\bookchapter{บทที่ 28:ใช้เงินอย่างคุ้มค่า ซื้อเครื่องที่ใช่—คู่มือเลือกซื้อโทรศัพท์ขั้นเทพ}

\chapterintro{เรื่องการซื้อโทรศัพท์นั้น หลายคนอ่านรีวิวมาสัปดาห์ก่อน เปรียบเทียบราคามาเป็นเดือน แต่สุดท้ายก็ยังเผลอจ่ายไปห้าถึงหกพันเพื่อซื้อเรือธงที่ตัวเองแทบไม่ได้ใช้ บทนี้จะไม่แสดงตารางสเปก เพราะสิ่งเหล่านั้นคุณดูได้จากเว็บไซต์อย่างเป็นทางการ สิ่งที่ผมอยากทำคือ มอบ“วิธีคิดเลือกซื้อมือถือ” ให้คุณ เพื่อให้ไม่ว่าตลาดจะออกโทรศัพท์รุ่นใหม่อะไรมา คุณก็สามารถตัดสินใจได้ด้วยตัวเองว่ามันเหมาะกับคุณหรือไม่}

\booksection{วิธีเลือกโทรศัพท์ 3 ขั้นตอน: ความต้องการ งบประมาณ รุ่น}

หลายคนถามบนอินเทอร์เน็ตว่า “สามพันซื้อมือถืออะไรดี” คำถามนี้ตอบไม่ได้หรอก เพราะไม่รู้ว่าคุณซื้อโทรศัพท์ไปใช้ทำอะไรเป็นหลัก ลำดับการคิดที่ถูกต้องในการซื้อโทรศัพท์ควรเป็น:

\begin{tcbinfo}{3 ขั้นตอนเลือกโทรศัพท์}
\begin{enumerate}
  \item \textbf{คิดให้ออกก่อนว่าคุณสนใจอะไรมากที่สุด}:ถ่ายภาพ? เล่นเกม? แบตเตอรี่? หน้าจอ? ระบบลื่นไหล? หรือแค่โทรและส่ง WeChat ได้ก็พอ? ให้จัดลำดับความต้องการ
  \item \textbf{แล้วก็กำหนดงบประมาณสูงสุด}:ไม่ใช่"ประมาณเท่าไหร่" แต่คือ"ยอมรับได้มากสุดเท่าไหร่" กำหนดขีดจำกัดที่ชัดเจนไว้ ตอนเลือกซื้อจะได้ไม่เลือกของที่แพงขึ้นเรื่อยๆ
  \item \textbf{ค่อยเลือกรุ่นเป็นขั้นสุดท้าย}:นำลำดับความต้องการและงบประมาณของคุณไปจับคู่กับรุ่นโทรศัพท์
\end{enumerate}

\end{tcbinfo}

\booksection{ดูสเปกหลักอย่างไร — ไม่ต้องเข้าใจหมด เข้าใจแค่ 5 ข้อนี้ก็พอ}

ตารางสเปกมีข้อมูลนับสิบรายการ แต่ผู้ใช้ทั่วไปแค่ใส่ใจ 5 มิตินี้ก็พอ:

\begin{center}
\begin{tabularx}{\linewidth}{|X|X|X|}
\hline
\textbf{มิติข้อมูล} & \textbf{วิธีดู} & \textbf{สรุปในหนึ่งประโยค} \tabularnewline
\hline
โปรเซสเซอร์ (SoC) & กำหนดว่าโทรศัพท์จะใช้งานลื่นไหลได้กี่ปี มาตรฐานปี 2026: Snapdragon 8 Elite Gen5 / Dimensity 9500 ขึ้นไป = ระดับท็อป; Snapdragon 8s Gen4 / Dimensity 8400 = ระดับกลาง-ระดับสูง; Snapdragon 7+ Gen4 / Dimensity 7350 = ใช้งานได้ 3 ปี ถ้าต่ำกว่าซีรีส์ Snapdragon 6 แนะนำให้พิจารณาให้ดี & ซื้อของใหม่ไม่ซื้อของเก่า SoC กำหนดว่าโทรศัพท์ของคุณจะใช้งานได้กี่ปีโดยไม่กระตุก \tabularnewline
หน้าจอ & OLED มีสีสันดีกว่าและประหยัดไฟกว่า LCD; อัตรารีเฟรช 120Hz เลื่อนจอลื่นไหลกว่า 60Hz เยอะมาก; ความละเอียดขั้นต่ำควรเป็น 1080P ขึ้นไป การปรับความสว่างแบบ PWM ความถี่สูง (1920Hz ขึ้นไป) จะดีต่อสายตามากกว่า & OLED + 120Hz + ปรับแสง PWM ความถี่สูง = มาตรฐานขั้นต่ำของหน้าจอที่ดี \tabularnewline
บลูทูธลำโพง & \textbf{ดูขนาดเซ็นเซอร์ ไม่ใช่จำนวนพิกเซล}เซ็นเซอร์ใหญ่กว่าย่อมดีกว่า ขนาดเซ็นเซอร์สำคัญกว่าจำนวนพิกเซลเยอะมาก กล้องเลนส์มุมกว้างพิเศษใช้งานได้จริง; กล้องซูม (ซูมออปติคอล 3x ขึ้นไป) สำหรับถ่ายวัตถุไกลๆ; กล้องมาโครส่วนใหญ่มีไว้ทำเลขเท่านั้น & เซ็นเซอร์ใหญ่ + ซูมใช้งานได้จริง > กล้องหลายตัวที่ทำเลข \tabularnewline
แบตเตอรี่และการชาร์จ & 4500mAh เป็นเกณฑ์ขั้นต่ำ, 5000mAh ขึ้นไปชาร์จวันละครั้งอุ่นใจได้เลย การชาร์จเร็วดูที่วัตต์: 60W ขึ้นไปถือว่าเร็ว, 100W ขึ้นไปเร็วมาก การชาร์จไร้สายสะดวกแต่ช้า & 5000mAh + ชาร์จไว 60W = คอมโบไร้กังวล \tabularnewline
การรับประกันอัปเดตระบบ & กำหนดว่าโทรศัพท์จะใช้งานได้อย่างปลอดภัยนานแค่ไหน Apple มักให้อัปเดตระบบ 5-6 ปี; แฟลกชิป Samsung ให้อัปเดตระบบ 4 ปี + อัปเดตความปลอดภัย 5 ปี; Google Pixel ให้ 7 ปี; สมาร์ทโฟนระดับกลางจากแบรนด์จีนมักให้การรับประกันอัปเดตสั้นกว่า (2-3 ปี) & อยากใช้นานๆ เลือกแบรนด์ที่รับประกันอัปเดตระบบนานๆ \tabularnewline
\hline
\end{tabularx}
\end{center}

\booksection{คำแนะนำแยกตามงบประมาณ (อ้างอิงตลาดกลางปี 2026)}

คำแนะนำต่อไปนี้อ้างอิงสภาพตลาดในจีนแผ่นดินใหญ่ปี 2026 ราคาเป็นราคาเริ่มต้นอย่างเป็นทางการ ราคาจริงอาจมีการปรับเปลี่ยน ขึ้นอยู่กับราคาจริงในขณะที่คุณซื้อ

\booksubsection{ต่ำกว่า 1000 หยวน: โทรได้ เล่น WeChat ได้ พอใช้ก็พอ}

\textbf{เหมาะกับใคร}นักเรียน, มือถือสำรอง, คนที่ไม่ได้เรียกร้องกับมือถือมาก, ใช้หลักๆ สำหรับโทรศัพท์และ WeChat

\textbf{ซื้ออะไรได้บ้าง}เวอร์ชันกลาง-ต่ำของ Redmi Note ซีรีส์, Honor X ซีรีส์, realme ซีรีส์ตัวเลข โปรเซสเซอร์ประมาณ Snapdragon ซีรีส์ 6 หรือ Dimensity 6000 ซีรีส์, หน้าจอเป็น LCD เป็นหลัก ถ่ายรูปได้ชัดแตะโฟกัสช้า, ความลื่นไหลของระบบถือว่าธรรมดา

\textbf{คำแนะนำในการเลือกซื้อ}ในราคานี้ให้ดู 2 อย่าง: แบตเตอรี่ใหญ่พอไหม (เลือก 5000mAh ขึ้นไป), แรมพอไหม (เลือกเริ่มต้น 6GB+128GB, อย่าซื้อแบบ 4GB เด็ดขาดจะกระตุก) Redmi มีข้อได้เปรียบชัดเจนที่สุดในราคานี้

\booksubsection{1000-2000 หยวน: คุ้มค่าที่สุด ใช้งาน 3-4 ปีไม่กระตุกในโซนราคาหวาน}

\textbf{เหมาะกับใคร}นักเรียน, งบจำกัดแต่อยากใช้นานๆ, ไม่ได้เรียกร้องสุดยอดเรื่องกล้องและเกม

\textbf{รุ่นที่เป็นตัวแทน}Redmi Turbo ซีรีส์, iQOO Z ซีรีส์, Honor X ซีรีส์เวอร์ชันสูง, OnePlus Ace ซีรีส์รุ่นเริ่มต้น

\textbf{คำแนะนำในการเลือกซื้อ}เลือกรุ่นที่มาพร้อมโปรเซสเซอร์ Snapdragon 7+ Gen3 หรือ Dimensity 7300 ขึ้นไป พร้อมแรม 8GB+256GB สเปกนี้ใช้งานประจำวัน + เล่น Honor of Kings เป็นบางครั้งได้สบายๆ ใช้งานได้ 3-4 ปีโดยไม่กระตุก

\booksubsection{2000-3500 หยวน: กำลังหลัก ทั้งประสิทธิภาพและกล้องถ่ายรูปดีขึ้นอย่างเห็นได้ชัด}

\textbf{เหมาะกับใคร}ผู้ใช้ทั่วไปที่อยากได้กล้องที่ไม่เลว, เล่นเกมไม่กระตุก, และหวังใช้งานได้นานๆ นี่คือช่วงราคาที่คนซื้อมากที่สุด

\textbf{รุ่นที่เป็นตัวแทน}Xiaomi ซีรีส์ตัวเลข (เช่น Xiaomi 17 เวอร์ชันมาตรฐานหลังลดราคา), vivo S ซีรีส์เวอร์ชันสูง, OPPO Reno ซีรีส์, Honor ซีรีส์ตัวเลข, Redmi K ซีรีส์

\textbf{คำแนะนำในการเลือกซื้อ}จุดเด่นที่สุดของช่วงราคานี้คือเป็น “มือถือรอบด้าน” — ไม่มีจุดอ่อนที่เห็นชัด เมื่อเทียบกับเครื่องราคาต่ำกว่า 2000 หยวน การอัปเกรดหลักคือ: กล้องถ่ายรูปดีขึ้นอย่างเห็นได้ชัด (มี OIS แล้ว), หน้าจอเป็น OLED คุณภาพสูง, ชาร์จไวขึ้น (พบเห็น 80W ขึ้นไปได้ทั่วไป), วัสดุและการประกอบดูพรีเมียมขึ้น เลือก 12GB+256GB หรือ 12GB+512GB ใช้งานได้ 4-5 ปีไม่มีปัญหา

\booksubsection{3500-5000 หยวน: ระดับก่อนแฟลกชิป สำหรับคนทั่วไปถือว่าเพียงพอแล้วยกเว้นนักเล่นฮาร์ดแวร์}

\textbf{เหมาะกับใคร}ผู้ใช้ที่มีความต้องการเรื่องกล้องถ่ายรูป และผู้ที่อยากสัมผัสความพรีเมียมระดับสูงสุด

\textbf{รุ่นที่เป็นตัวแทน}Xiaomi ซีรีส์ตัวเลขเวอร์ชันมาตรฐาน, vivo X ซีรีส์, OPPO Find X ซีรีส์เวอร์ชันมาตรฐาน, แฟลกชิปเก่าของ iPhone (เช่น iPhone 16)

\textbf{คำแนะนำในการเลือกซื้อ}ในราคานี้คุณสามารถซื้อแฟลกชิปของปีที่แล้วได้แล้ว คำแนะนำของผมคือ:\textbf{ซื้อแฟลกชิปปีที่แล้ว ดีกว่ามือถือระดับกลางของปีนี้}วัสดุ การประกอบ และฮาร์ดแวร์กล้องของแฟลกชิปปีที่แล้วมักจะดีกว่าระดับกลางของปีนี้ไปหนึ่งระดับ ทั้งนี้เงื่อนไขคือคุณต้องยอมรับว่ามันไม่ใช่ “รุ่นล่าสุด”

\booksubsection{มากกว่า 5000 หยวน: โซนแฟลกชิป สำหรับคนไม่ขาดเงินหรือผู้ที่ไล่ล่าประสบการณ์สูงสุดในบางด้าน}

\textbf{เหมาะกับใคร}สายถ่ายภาพ, นักเล่นเกมหนัก, ผู้ที่ต้องการหน้าจอและงานประกอบระดับท็อป, คอนเทนต์ครีเอเตอร์ที่ใช้มือถือถ่ายวิดีโอ

\textbf{รุ่นที่เป็นตัวแทน}iPhone 17 Pro/Pro Max, Samsung S26 Ultra, Huawei Pura 80 Ultra, Xiaomi 17 Ultra, vivo X300 Pro

\textbf{คำแนะนำในการเลือกซื้อ}ในระดับราคานี้ความแตกต่างระหว่างแต่ละแบรนด์น้อยลงเรื่อยๆ จริงๆ แล้วขึ้นอยู่กับว่าคุณชอบระบบปฏิบัติการอะไร ชอบถ่ายรูปเลือก vivo หรือ Huawei, ชอบถ่ายวิดีโอเลือก iPhone, ชอบหน้าจอและ S Pen เลือก Samsung, ชอบความสมดุลเลือก Xiaomi พูดตรงๆ ถึงระดับราคานี้แล้ว,\textbf{คุณเลือกตามใจตัวเองได้เลย แบบไหนสบายตาก็เลือกเครื่องนั้น}

\booksection{เลือกมือถือตามการใช้งาน — คุณต้องการอะไร?}

ด้านล่างนี้คือ “โปรไฟล์ผู้ใช้งาน” ที่พบได้บ่อยที่สุด หาที่ตรงกับตัวเองได้เลย:

\begin{center}
\begin{tabularx}{\linewidth}{|X|X|X|}
\hline
\textbf{สถานการณ์การใช้งาน} & \textbf{สิ่งที่ควรดูเป็นหลัก} & \textbf{ทิศทางที่แนะนำ} \tabularnewline
\hline
สายถ่ายภาพ & เซ็นเซอร์หลักขนาดใหญ่ (1 นิ้วดีที่สุด), กันสั่นออปติคอล, ซูม (ซูมออปติคอล 3x ขึ้นไป), ชิปประมวลผลภาพ & vivo X ซีรีส์, Huawei Pura/Mate ซีรีส์, Xiaomi Ultra ซีรีส์ \tabularnewline
สายเกมเมอร์ & โปรเซสเซอร์ระดับท็อป, อัตราการสุ่มสัมผัสสูง, ระบบระบายความร้อน, แบตเตอรี่ใหญ่ + ชาร์จไว & iQOO ซีรีส์ตัวเลข, ROG Phone, Red Magic, Redmi K ซีรีส์ \tabularnewline
เน้นแบตเตอรี่ & แบตเตอรี่ 5500mAh ขึ้นไป, โปรเซสเซอร์ประหยัดพลังงาน, หน้าจอ LTPO & Honor Magic Series, OnePlus, Redmi Note เวอร์ชันแบตใหญ่ \tabularnewline
ผู้สูงอายุใช้งาน & โหมดใช้งานง่าย, ฟอนต์ใหญ่, เสียงดัง, แบตเตอรี่ใหญ่, มีจุดบริการหลังการขายมาก & Huawei Enjoy, Honor Play, OPPO A Series, Redmi Note \tabularnewline
นักเรียน/นักศึกษา & คุ้มค่า, ทนทาน, แบตเตอรี่ใหญ่, มีฟังก์ชันควบคุมโดยผู้ปกครอง & Redmi Note/Turbo, realme GT, iQOO Neo \tabularnewline
ผู้ใช้ในระบบนิเวศ Apple & ถ้ามี Mac/iPad/AirPods อยู่แล้ว การซื้อ iPhone คือคำตอบเดียว & iPhone 17 เวอร์ชันมาตรฐาน (คุ้มค่า) หรือเวอร์ชัน Pro (ประสบการณ์เต็มรูปแบบ) \tabularnewline
เดินทางไปต่างประเทศ & รองรับ eSIM, รองรับคลื่นความถี่ทั่วโลกครบถ้วน, มีบริการ Google & iPhone, Samsung S Series, Pixel, OnePlus \tabularnewline
\hline
\end{tabularx}
\end{center}

\booksection{ห้าข้อผิดพลาดที่มักเกิดขึ้น — อ่านให้แน่ใจก่อนซื้อ}

\begin{enumerate}
  \item \textbf{ข้อผิดพลาดที่ 1: กล้องสี่หรือห้าตัวไม่ได้แปลว่าถ่ายรูปดี}— แบรนด์มักชอบเพิ่มจำนวนเลนส์ในสเปก ซึ่งมักจะมีอยู่ 2 ตัวที่เป็นเลนส์ 200 ล้านพิกเซล "เติมจำนวน" (ถ่ายมาโคร + วัดความลึก) ซึ่งจริงๆ แล้วใช้ประโยชน์ไม่ได้เลย ถ้าจะดูความสามารถในการถ่ายรูปของมือถือ\textbf{ให้ดูแค่กล้องหลัก + อัลตร้าไวด์ + เทเลโฟโต้}สามตัวนี้ก็เพียงพอแล้ว ส่วนที่เกินมาถือว่าไม่มีอยู่จริงก็ได้
  \item \textbf{ข้อผิดพลาดที่ 2: อย่าถูกหลอกด้วย 100 ล้านพิกเซล}— พิกเซลสูงไม่ได้แปลว่าคุณภาพภาพดี ขนาดเซ็นเซอร์, รูรับแสง, การปรับแต่งอัลกอริทึม สำคัญกว่าพิกเซลมาก มือถือแฟลกชิปหลายรุ่นใช้เซ็นเซอร์ขนาดใหญ่ 50 ล้านพิกเซล ถ่ายภาพออกมาดีกว่ามือถือราคาพันกว่าบาทที่มี 100 ล้านพิกเซลมาก
  \item \textbf{ข้อผิดพลาดที่ 3: อย่าซื้อมือถือเครื่องผูกสัญญาเครือข่าย}— "ฝากเครดิตแล้วแถมมือถือ" ดูเหมือนจะคุ้มมาก แต่มือถือผูกสัญญาเครือข่ายมักจะอัปเดตระบบช้า, อาจมีแอปแถมที่ลบไม่ได้, และมักจะถูกล็อกให้ใช้เครือข่ายเดียว (ใช้ซิมเครือข่ายนี้ได้เท่านั้น) ดูเหมือนจะประหยัดเงิน แต่จริงๆ แล้วคือการ "ล็อก" คุณไว้
  \item \textbf{ข้อผิดพลาดที่ 4: "หน้าจอโค้งดูพรีเมียม" — นั่นคือการตลาด ไม่ใช่ความจริงเสมอไป}— หน้าจอโค้งเคยเป็นสัญลักษณ์ของมือถือแฟลกชิป แต่ตอนนี้คนจำนวนมากขึ้นพบปัญหาของมัน: แปะฟิล์มยาก, แตะผิดพลาดที่ขอบ, และถ้าแตกซ่อมแพง แนวโน้มในปี 2026 คือการหันกลับไปใช้จอ "โค้งนิดหน่อย" หรือแม้กระทั่ง "จอแบนราบ" อย่าซื้อจอโค้งเพียงเพราะคิดว่า "โค้งแล้วดูพรีเมียม"
  \item \textbf{ข้อผิดพลาดที่ 5: ไม่ใช่มือถือ "AI" ทุกรุ่นที่จะมีประโยชน์จริง}— ในปี 2026 ทุกแบรนด์ต่างโหระพร "มือถือ AI" แต่หลายรุ่นแค่เปลี่ยนชื่อผู้ช่วยเสียงที่มีอยู่แล้วเท่านั้น ฟังก์ชัน AI ที่มีประโยชน์จริงๆ คือ: การแปลภาษาแบบเรียลไทม์, การลบภาพด้วย AI (ลบคนออกจากรูป), แต่งรูปด้วย AI, การประมวลผลโมเดลภาษาขนาดใหญ่บนเครื่อง, และสรุปเนื้อหาอัจฉริยะแบบใช้งานในพื้นที่ ฟังก์ชันอื่นๆ ที่ดูโหวงๆ ส่วนใหญ่เป็นแค่เครื่องมือหลอกล่อ
\end{enumerate}

\booksection{เมื่อไหร่ที่ซื้อจะคุ้มที่สุด?}

\begin{itemize}
  \item \textbf{618 (ปลายเดือนพฤษภาคมถึงกลางเดือนมิถุนายน)}: ช่วงเวลาที่ราคามือถือลดลงมากที่สุดในรอบปี ผู้ผลิตมือถือหลายรายต้องการเพิ่มยอดขาย จึงลดราคามือถือแฟลกชิปได้ถึง 800-1500 หยวน เป็นช่วงเวลาที่ดีที่สุดในการซื้อมือถือในครึ่งปีแรก
  \item \textbf{Double 11 (ปลายเดือนตุลาคมถึงกลางเดือนพฤศจิกายน)}: โปรโมชันใหญ่ในครึ่งปีหลัง ความรุนแรงของส่วนลดใกล้เคียงกับ 618 ในช่วงเวลานี้แบรนด์หลายเจ้าเพิ่งเปิดตัวมือถือแฟลกชิปรุ่นใหม่ ทำให้มือถือแฟลกชิปรุ่นเก่าลดราคาลงมากที่สุด
  \item \textbf{หนึ่งเดือนหลังเปิดตัวเครื่องใหม่}: ตอนเปิดตัวใหม่ๆ จะขายในราคาปกติ หลังจากนั้นหนึ่งเดือน ราคาในช่องทางจำหน่าย (ร้านค้าบุคคลที่สาม) มักจะถูกลง 200-500 หยวน
  \item \textbf{สองสัปดาห์ก่อนตรุษจีนถึงเทศกาลโยวเซียว}: เทศกาลสินค้าตรุษจีนเพื่อเคลียร์สต็อกของร้านค้าออฟไลน์ สามารถต่อรองราคากับพนักงานในร้านได้
  \item \textbf{ช่วงเวลาที่ไม่แนะนำ}: ซื้อเป็นคนแรกๆ ในวันที่จัดงานเปิดตัวสินค้าใหม่ — ยกเว้นคุณจะอยากเป็นกลุ่มคนแรกที่ได้ใช้จริงๆ ไม่งั้นราคาเปิดตัวจะแพงที่สุด
\end{itemize}

\booksection{เลือกช่องทางไหนดี — ซื้อที่ไหนถึงน่าเชื่อถือ?}

\begin{itemize}
  \item \textbf{แบรนด์แมกกาซีนอย่างเป็นทางการ / แฟลกชิปสโตร์ Tmall / JD สินค้าแท้}: ปลอดภัยที่สุด ราคาอาจไม่ดีนัก แต่ไม่ต้องกังวลเรื่องของปลอมหรือของมือสองทำใหม่ และบริการหลังการขายดีที่สุด
  \item \textbf{Pinduoduo "ส่วนลดหมื่นล้าน"}: ราคาถูกที่สุด แต่ต้องเลือกร้านที่ระบุว่าเป็น "ส่วนลดหมื่นล้าน" ร้านที่ไม่ระบุอย่าไปยุ่งกับมัน เมื่อได้รับสินค้าแล้วอย่าลืมถ่ายคลิปแกะกล่อง ถ้ามีปัญหาจะได้คืนของได้
  \item \textbf{ร้านค้าออฟไลน์}: สามารถจับและลองสัมผัสจริงได้ ข้อควรระวัง: พนักงานขายจะแนะนำรุ่นที่กำไรสูง (แทนที่จะเป็นรุ่นที่เหมาะกับคุณที่สุด) แนะนำให้เลือกเป้าหมาย 2-3 รุ่นทางออนไลน์ก่อนเข้าร้าน พอไปถึงก็บอกไปเลยว่า "อยากดูแค่รุ่นเหล่านี้"
  \item \textbf{แพลตฟอร์มมือสอง (Xianyu / Zhuanzhuan)}: ถูกกว่าแต่มีความเสี่ยง ถ้าจะซื้อของมือสอง ให้จำกฎสามข้อไว้: ขอใบเสร็จหรือสกรีนช็อตคำสั่งซื้อ, ขอหมายเลข IMEI เพื่อตรวจสอบ, และ\textbf{ทำธุรกรรมผ่านแพลตฟอร์มเท่านั้น ห้ามโอนเงินส่วนตัวเด็ดขาด}
\end{itemize}

\begin{tcbpractice}{ลงมือทำ: ใช้วิธีนี้เลือกมือถือสักเครื่อง}
\begin{enumerate}
  \item หยิบกระดาษมาแผ่น แล้วเขียนความต้องการหลัก 2 อย่างของคุณลงไป (เช่น: แบตเตอรี่เป็นอันดับหนึ่ง, กล้องเป็นอันดับสอง)
  \item กำหนดงบประมาณสูงสุดไว้ (เช่น 3500 หยวน)
  \item เปิด JD หรือ Tmall กรองตามราคา แล้วเลือกเครื่องเปรียบเทียบไว้ 3-5 รุ่น
  \item ไปค้นหาใน Bilibili ว่า "รุ่น XX ใช้งานจริง 1 เดือน" (ไม่ใช่รีวิววันแรก แต่เป็นความรู้สึกจริงหลังใช้งานอย่างน้อยหนึ่งเดือน) แล้วดูสัก 2-3 คลิป
  \item ยืนยันว่าช่องทางการซื้อน่าเชื่อถือแล้วสั่งซื้อ และอัดวิดีโอแกะกล่องตอนรับสินค้า
\end{enumerate}

\end{tcbpractice}

\begin{tcbsummary}{สรุปบทนี้}
\begin{itemize}
  \item 3 ขั้นตอนเลือกสมาร์ทโฟน: คิดความต้องการก่อน กำหนดงบประมาณ เลือกรุ่น
  \item สเปกหลักดูแค่ 5 อย่าง: โปรเซสเซอร์, หน้าจอ, กล้อง, แบตเตอรี่/การชาร์จ, และการรับประกันอัปเดตระบบ
  \item ช่วงราคา 1500-3500 หยวน คือจุดที่คุ้มค่าที่สุดและแนะนำมากที่สุด
  \item ซื้อแฟลกชิปปีที่แล้ว คุ้มกว่าซื้อมิดเรนจ์ปีนี้
  \item 5 กับดัก: กล้องหลายตัวที่หลอก, กับดักพิกเซลสูง, เครื่องคอมกับโปรแกรมเครือข่าย, ความเชื่อเรื่องจอโค้ง, และ AI ที่เป็นแค่เทคนิคการตลาด
  \item 618 และ Double 11 คือ 2 ช่วงเวลาที่เหมาะที่สุดในการซื้อมือถือในรอบปี
\end{itemize}

\end{tcbsummary}
